\section{A simplex gradient is a multidirectional finite difference}
\subsection{Choose the metric before choosing the neighborhood}
Distances between coordinates with different units have no canonical meaning.
Even when all coordinates are currents, different allowed ranges can dominate
the triangulation. Declare reference scales $q_j>0$, let
$\matr Q=\diag(q_1,\ldots,q_d)$, and use
\begin{equation}
 \vect u=\matr Q^{-1}(\vect x-\vect x_*).
 \label{eq:metric}
\end{equation}
All examples use dimensionless coordinates; the code triangulates exactly
the coordinates supplied by its caller. Automatic rescaling would hide the
choice of locality. Translation and common scaling preserve the Euclidean
geometry in exact arithmetic; anisotropic scaling can change the neighbors.

If $\widetilde f(\vect u)=f(\vect x_*+\matr Q\vect u)$, the chain rule gives
\begin{equation}
 \nabla_{\vect x} f=\matr Q^{-\mathsf T}\nabla_{\vect u}\widetilde f,
 \qquad
 \matr H_{\vect x}=\matr Q^{-\mathsf T}\matr H_{\vect u}\matr Q^{-1}.
 \label{eq:chain}
\end{equation}
A gradient component has field units divided by its coordinate unit;
Hessian entries divide by two coordinate units. A normalized gradient step
in $\vect u$ selects a search metric and generally differs from a Euclidean
steepest step in physical $\vect x$ coordinates.

\subsection{Subtracting the reference vertex exposes the slope}
Consider a nondegenerate simplex $\sigma$ with vertices
$\vect x_0,\ldots,\vect x_d$. Define
$\vect e_j=\vect x_j-\vect x_0$ and
\begin{equation}
 \matr E_\sigma=
 \begin{bmatrix}\vect e_1^{\mathsf T}\\\vdots\\\vect e_d^{\mathsf T}\end{bmatrix},
 \qquad
 \Delta\vect f_\sigma=
 \begin{bmatrix}f_1-f_0\\\vdots\\f_d-f_0\end{bmatrix}.
\end{equation}
\glsadd{sym:simplexedge}\glsadd{sym:recoveredgradient}
An affine interpolant is
\begin{equation}
 \widehat f_\sigma(\vect x)=f_0+
 \vect g_\sigma^{\mathsf T}(\vect x-\vect x_0).
\end{equation}
Matching the other vertices gives the central solve
\begin{equation}
 \boxed{\matr E_\sigma\vect g_\sigma=\Delta\vect f_\sigma.}
 \label{eq:edgesolve}
\end{equation}
Each row specifies one directional change. Independent edges determine all
gradient components. A nearly collapsed simplex leaves a direction poorly
observed, so small value changes can require large slopes.

The equivalent augmented solve uses rows $[\vect x_j^{\mathsf T},1]$ and
solves for slope and intercept simultaneously. Subtraction eliminates the
intercept first and avoids a large coordinate offset in that matrix.
Barycentric interpolation is another view: for affine barycentric weights
$\lambda_j$, $\widehat f=\sum_j\lambda_j f_j$ and
$\vect g=\sum_j f_j\nabla\lambda_j$. These describe the same $P_1$
interpolant. A solve through exactly $d+1$ values neither averages nor
regularizes measurement noise.

For $f(\vect x)=\vect a^{\mathsf T}\vect x+b$,
$\Delta\vect f_\sigma=\matr E_\sigma\vect a$, hence
$\vect g_\sigma=\vect a$ on every admissible cell. This affine reproduction
is an algebraic invariant, not an empirical accuracy claim.

\subsection{A normalized triangle makes the approximation visible}
Take $(0,0)$, $(1,0)$ and $(0,1)$. If their values are $0$, $2$ and
$-1$, the edge matrix is the identity and the gradient is $(2,-1)$.
The plane increases by two units in the first coordinate and decreases by
one in the second. Its level lines are perpendicular to that vector.

For $f=(x_1^2+x_2^2)/2$ on the same vertices, the affine gradient is
$(1/2,1/2)$, while the exact gradient at centroid $(1/3,1/3)$ is
$(1/3,1/3)$. Interpolating the values exactly does not differentiate the
underlying nonlinear field exactly at the representative. The residual at
every vertex is zero in both examples, so interpolation residual cannot
diagnose gradient error.

\begin{figure}[htbp]
\centering\begin{tikzpicture}
\begin{scope}[scale=2.4]
\draw[constraint curve] (0,0)--(1,0)--(0,1)--cycle;
\draw[coordinate vector] (0,0)--(1.35,0) node[annotation,right] {$x_1$};
\draw[coordinate vector] (0,0)--(0,1.2) node[annotation,left] {$x_2$};
\node[annotation,below left] at (0,0) {$f_0=0$};
\node[annotation,below] at (1,0) {$f_1=2$};
\node[annotation,left] at (0,1) {$f_2=-1$};
\draw[eigenvector] (.33,.33)--(.93,.03);
\node[annotation,above right] at (.5,.5) {$\vect g=(2,-1)$};
\end{scope}
\begin{scope}[xshift=6cm,scale=2.4]
\draw[constraint curve] (0,0)--(1,0)--(.4,.12)--cycle;
\draw[coordinate vector] (0,0)--(1.2,0) node[annotation,below] {$x_1$};
\draw[coordinate vector] (0,0)--(0,.8) node[annotation,left] {$x_2$};
\draw[projection line] (.4,0)--(.4,.12);
\node[annotation,above] at (.4,.12) {$(a,\delta)$};
\node[annotation,above right] at (.2,.62) {$g_2=(\Delta f_2-a\Delta f_1)/\delta$};
\end{scope}
\end{tikzpicture}

\caption{An affine slope is recovered exactly, but nonlinear slopes are
approximated. A thin triangle magnifies vertical value error: the second
slope contains division by the height $\delta$.}
\label{fig:simplex}
\end{figure}

\subsection{Noise amplification and curvature bias are different errors}
Write observations as $f_j=f(\vect x_j)+\eta_j$. Then
\begin{equation}
 \delta\vect g_\sigma=\matr E_\sigma^{-1}\Delta\vect\eta_\sigma,
 \qquad
 \norm{\delta\vect g_\sigma}_2\leq
 \norm{\matr E_\sigma^{-1}}_2\norm{\Delta\vect\eta_\sigma}_2.
 \label{eq:noisebound}
\end{equation}
For independent errors of common variance $\sigma_\eta^2$, every difference
shares $\eta_0$. Their covariance is therefore not diagonal:
\begin{align}
 \operatorname{Cov}(\Delta\vect\eta)&=
 \sigma_\eta^2(\matr I+\vect 1\vect 1^{\mathsf T}),\\
 \operatorname{Cov}(\delta\vect g)&=\sigma_\eta^2
 \matr E^{-1}(\matr I+\vect1\vect1^{\mathsf T})\matr E^{-\mathsf T}.
 \label{eq:covariance}
\end{align}
For unequal or correlated errors, replace the middle covariance by the one
induced by the differencing operator. The implementation records
$\norm{\matr E^{-1}}_2$ and edge condition number separately. A tiny
equilateral cell can be well-conditioned and still amplify noise strongly.
Filtering only by condition number cannot address that situation.

If the Hessian norm of a twice differentiable $f$ is bounded by $L$ on the
cell, Taylor expansion at $\vect x_0$ gives remainders
$|r_j|\leq L\norm{\vect e_j}^2/2$. Thus
\begin{align}
 \norm{\vect g_\sigma-\nabla f(\vect x_0)}_2
 &\leq \frac{L}{2}\norm{\matr E_\sigma^{-1}}_2
 \left(\sum_{j=1}^d\norm{\vect e_j}_2^4\right)^{1/2}\nonumber\\
 &\quad+\norm{\matr E_\sigma^{-1}}_2\norm{\Delta\vect\eta}_2.
 \label{eq:bias}
\end{align}
Moving the reference to the centroid adds at most
$L\norm{\vect z_\sigma-\vect x_0}$ to the bound. On shape-regular cells
of characteristic length $\rho$, curvature bias decreases like $L\rho$,
whereas fixed-amplitude noise is amplified like $\sigma_\eta/\rho$.
More samples do not automatically mean better derivatives.

The determinant measures oriented volume up to the factor $d!$. Its
vanishing signals collapse, but its magnitude mixes size and shape.
Singular values separate poor directional support from overall scale.
Delaunay connectivity does not ensure well-conditioned cells, especially
in higher dimensions or near support boundaries.
